\documentclass{memoir}
\usepackage[utf8]{inputenc}
\usepackage[T1]{fontenc}
\usepackage{geometry}
\usepackage{graphicx}
\usepackage{float}
\usepackage{hyperref}
\usepackage[english,indonesian]{babel}

\title{Optimasi Pemrosesan Template LaTeX Menggunakan Pendekatan Berbasis AI}
\author{Academic Support Agent}
\date{\today}

\begin{document}
\maketitle

\begin{abstract}

\end{abstract}

\chapter{PENDAHULUAN}


\chapter{TINJAUAN PUSTAKA}


% Original content (commented out):
% \documentclass[a4paper,oneside,article]{memoir}
\usepackage[utf8]{inputenc}  % Korrekt håndtering af æ, ø og å
\usepackage[T1]{fontenc}  % Korrekt håndtering af æ, ø og å
\usepackage{microtype}  % Typografisk magi! Giver bl.a. pænere orddeling
\usepackage{graphicx}  % Gør det muligt at indsætte billeder
\usepackage{amsmath}  % Giver adgang til uundværlige matematikting
\usepackage[danish]{babel}  % Danske betegnelser og orddeling
\renewcommand{\danishhyphenmins}{22}  % Bedre dansk orddeling

% Indholdet i dit dokument starter her!
\begin{document}
\author{Fornavn Efternavn -- Studienummer}
\title{Grundskabelon til MFSR}
\date{\today}
\maketitle   % Indsætter overskriften

Velkommen til AU, naturvidenskab!

\end{document}
% Dokumentet er nu slut! Skriv ikke noget herunder.

\end{document}